\section{The curvature covariance on the selected tree slice}
\label{sec:actual-reader-tree-covariance}

The local Gaussian comparison needs an exact identification at the flat
configuration. Its reference matrix is the inverse on the selected forest,
whereas the free curvature estimates are expressed using two consecutive
members of the precision tower. We now identify these two descriptions.
Everything in this section is on the same pair of finite tori. In
particular, no limit in the period and no change of forest is involved.

\subsection{Counting coordinates and the flat precision}

Write $E_f^r$ and $E_c^r$ for real $r$-cochains on the fine and coarse
tori, with their counting inner products. Set
\[
 d_f:E_f^0\longrightarrow E_f^1,\quad
 D:E_f^1\longrightarrow E_f^2,\quad
 d_c:E_c^0\longrightarrow E_c^1,\quad
 D':E_c^1\longrightarrow E_c^2.
\]
Thus $Dd_f=0$ and $D'd_c=0$. Matrices initially act on one colour;
tensoring with the identity on the real orthonormal colour space gives
the physical arrays. There is no volume factor in this convention.

Let $M_f,M_c$ be the actual symmetric two-form metrics at the two levels,
with $M_f\ge aI$ and $M_c\ge a'I$ for $a,a'>0$. Put
\[
 K=D^{\mathsf T}M_fD,\qquad
 K'=(D')^{\mathsf T}M_cD',\qquad
 Q=\overline Q+d_cR_0,\qquad
 A=K+\kappa Q^{\mathsf T}Q,\quad \kappa>0.
\]
Here $\overline Q:E_f^1\to E_c^1$ is the selected averaged path
map, $R_0:E_f^1\to E_c^0$ is its root correction, and $Q$ is the
derivative of the actual predictor. If $E:E_f^0\to E_c^0$ evaluates
a potential at the block roots, then
\begin{equation}\label{eq:reader-root-gradient}
 Qd_f=d_cE.
\end{equation}
Indeed the integral of a gradient along a path is the difference of
its endpoint values. In $\overline Qd_f$ these are the weighted
endpoint averages. The root correction replaces each such average by
the value at the root, giving \eqref{eq:reader-root-gradient}.
This also explains why a gradient with zero root values costs no
energy in $A$.

We first relate $A$ to the flat rows used in the local comparison.
Suppose the selected path coefficients are absolutely summable and
attached to a two-form row $R_f$: the sum of the coefficients with
each given endpoint and plaquette type is exactly that row entry.
Suppose the full factorization and symmetry identities are
\begin{equation}\label{eq:reader-flat-factor}
 M_f=\alpha I+R_f^{\mathsf T}R_f,\qquad \alpha>0,
\end{equation}
and the curl energy $u^{\mathsf T}Kv$ is invariant under the selected
signed lattice motions. These are the path attachment, metric
factorization and symmetry properties of the same fine bed; they
must refer to that bed, not an independently chosen comparison metric.

At the zero background, with every path admitted, parallel transport
in each path row is the identity. The coefficient sum in that row is
therefore $R_f$, by absolute summability and attachment. The fine
linear response in channel $s$ is $DP_s$, where $P_s$ is the signed
permutation on link cochains. The predictor response is the signed
permutation of $Qu$ on coarse bonds. Consequently the channel's
bilinear form is
\[
 \alpha\langle DP_su,DP_sv\rangle
 +\langle R_fDP_su,R_fDP_sv\rangle
 +\kappa\langle Qu,Qv\rangle
 =\langle u,Av\rangle.
\]
For the equality, first use \eqref{eq:reader-flat-factor}, then the
stated invariance of the curl energy; the coarse signed permutation
preserves the counting inner product. Each channel has this same
bilinear form, so its normalized orientation average is also $A$.
Restricting to the selected tree gives the tree compression of $A$.
This identification is at zero background with all paths admitted;
the estimates comparing a curved or cut row with it are separate.

\subsection{Rooted gauge projection and the slice inverse}

Fix the selected forest $F$, with one designated root in each block.
Let $I_F$ insert the coordinates of the nonforest bonds into $E_f^1$,
putting zero on every forest bond. The coordinate space of these
free bonds will be denoted $V_F$. For $u\in E_f^1$ define $\phi_u(x)$
as the oriented integral of $u$ along the forest path from $x$ back
to its root. Define
\[
 \mathcal Pu=u+d_f\phi_u.
\]
The root values of $\phi_u$ vanish. The difference of two adjacent
path integrals is the negative of the intervening oriented bond
value, so $\mathcal Pu$ vanishes on the forest. If $u$ already
vanishes there, all its forest path integrals vanish, and
$\mathcal Pu=u$. Furthermore,
\begin{equation}\label{eq:reader-slice-projection}
 A\mathcal P=A,\qquad D\mathcal P=D.
\end{equation}
The second equality follows from $Dd_f=0$; the first follows from
the same identity and \eqref{eq:reader-root-gradient}, since
$E\phi_u=0$.

The tree comparison of Section~\ref{sec:sector-gram}, specialized
to the flat field and summed over bonds, gives finite constants
$b_1,b_2\ge0$ such that
\[
 \|u\|^2\le b_1\|Du\|^2+b_2\|Qu\|^2
 \quad\text{whenever }u\text{ vanishes on }F.
\]
Here the block geometry and predictor branch hypotheses of that
comparison are retained. In particular the fine side is at least
one hundred times the macro radius and the coarse side exceeds one,
as in its finite-torus application. Hence the compressed precision
$A_F=I_F^{\mathsf T}AI_F$ is positive definite, with lower bound
$\min(a,\kappa)/(1+b_1+b_2)$. Define its extended covariance by
\[
 X=I_FA_F^{-1}I_F^{\mathsf T}.
\]
For any $u$, write $\mathcal Pu=I_Fv$ in its free coordinates.
Using \eqref{eq:reader-slice-projection} gives
\[
 XAu=XA\mathcal Pu
 =I_FA_F^{-1}I_F^{\mathsf T}AI_Fv=I_Fv=\mathcal Pu.
\]
We have proved the exact identity
\begin{equation}\label{eq:reader-curl-through}
 DXA=D.
\end{equation}
The full precision need not be invertible. Its inverse on the
forest slice is enough, because curl is unchanged by the rooted
gauge projection.

\subsection{Curvature matrices without inverting zero modes}

For any symmetric positive definite matrix $M$ on two-forms and
any finite curl matrix $q$, there is a matrix $Z$ satisfying
\begin{equation}\label{eq:reader-curl-solve}
 (q^{\mathsf T}Mq)Z=q^{\mathsf T}.
\end{equation}
To see this, positivity gives
$\ker(q^{\mathsf T}Mq)=\ker q$. Since $q^{\mathsf T}Mq$ is
symmetric, its range is the orthogonal complement of this kernel,
which is precisely $\operatorname{ran}q^{\mathsf T}$. Solve each
column of \eqref{eq:reader-curl-solve} in that range.
Any two solutions differ by a matrix with columns in $\ker q$.
Thus $P_M=qZ$ is independent of the choice. Moreover
\[
 P_M=Z^{\mathsf T}(q^{\mathsf T}Mq)Z,
\]
by transposing \eqref{eq:reader-curl-solve}. It is therefore
symmetric and positive semidefinite.

Apply this construction on both tori, choosing
\[
 KZ=D^{\mathsf T},\qquad K'Z'=(D')^{\mathsf T},\qquad
 P_f=DZ,\quad P_c=D'Z'.
\]
These are the finite-period free curvature matrices. This
construction does not divide by the curl eigenvalue at zero
momentum or discard a harmonic link direction.

Let $A_0:E_c^1\to E_f^1$ be the selected one-form transport.
The stationarity identities on the same beds are
\begin{equation}\label{eq:reader-stationarity-input}
 KA_0=\overline Q^{\mathsf T}K',\qquad
 \overline QA_0=I-\kappa^{-1}K'.
\end{equation}
For completeness, the two-form transport used below is
$J=H_f^{-1}T H_c$, where $H_f,H_c$ are the diagonal two-form
operators of the selected beds and $T$ is their installed tower
map. Their inverse and transport identities give
\begin{equation}\label{eq:reader-curvature-intertwiner}
 DA_0=JD'.
\end{equation}
The operators $H_f,H_c$ in this formula are the diagonal operators;
the metrics defining $K,K'$ remain $M_f,M_c$. This distinction is
needed when the metric has undergone its prescribed repair.

Straighten $A_0$ onto the forest by setting $W=A_0+d_f\Phi$, where
the column potential $\Phi$ is the backwards forest integral of
$A_0$, minus the blockwise constant root correction $R_0A_0$.
It has root values $E\Phi=-R_0A_0$, and the same path difference
calculation makes $W$ vanish on the forest. Equations
\eqref{eq:reader-root-gradient}--\eqref{eq:reader-stationarity-input}
give
\begin{equation}\label{eq:reader-stationary-slice}
 QW=I-\kappa^{-1}K',\qquad KW=Q^{\mathsf T}K',
 \qquad AW=\kappa Q^{\mathsf T},\qquad DW=JD'.
\end{equation}
Indeed $QW=\overline QA_0$ because the root correction cancels.
Also $KW=KA_0$, while
$(Q-\overline Q)^{\mathsf T}K'=R_0^{\mathsf T}d_c^{\mathsf T}K'=0$.
Adding $\kappa Q^{\mathsf T}QW$ to $KW$ gives the third identity.
The last follows from \eqref{eq:reader-curvature-intertwiner} and
$Dd_f=0$.

\subsection{Exact conditional curvature covariance}

The cross-covariance identity needed in the calculation is
\begin{equation}\label{eq:reader-cross-covariance}
 D'QZ=P_cJ^{\mathsf T}.
\end{equation}
Transpose $KW=Q^{\mathsf T}K'$ and multiply by $Z$. This yields
\[
 K'QZ=W^{\mathsf T}KZ=W^{\mathsf T}D^{\mathsf T}
       =(D')^{\mathsf T}J^{\mathsf T}.
\]
Since $(Z')^{\mathsf T}K'=D'$ and
$(Z')^{\mathsf T}(D')^{\mathsf T}=(D'Z')^{\mathsf T}=P_c$,
multiplication by $(Z')^{\mathsf T}$ proves
\eqref{eq:reader-cross-covariance}.

We can now compute directly. Equations
\eqref{eq:reader-curl-through} and \eqref{eq:reader-stationary-slice}
give $DX(\kappa Q^{\mathsf T})=DW$. Also
\[
 AZ=D^{\mathsf T}+\kappa Q^{\mathsf T}QZ.
\]
Multiplication by $DX$ and use of \eqref{eq:reader-curl-through}
therefore show
\[
 DXD^{\mathsf T}+DWQZ=DZ=P_f.
\]
Substitute $DW=JD'$ and
\eqref{eq:reader-cross-covariance}. We obtain
\begin{equation}\label{eq:reader-tree-conditional-covariance}
 (DI_F)A_F^{-1}(DI_F)^{\mathsf T}
       =P_f-JP_cJ^{\mathsf T}.
\end{equation}
In particular the right side is positive semidefinite, being the
covariance of the curl of a centered Gaussian on $V_F$.
Tensoring every map with the colour identity proves the same
identity for all physical colour coordinates and all cross-entries.

Equation~\eqref{eq:reader-tree-conditional-covariance} identifies the
flat reference covariance on the actual tree with the tower's
conditional curvature matrix. Together with the flat row
identification above, it supplies the exact algebraic reference
for the local Gaussian comparison. Passing from a posed, truncated
local covariance to this reference still uses the local inverse and
window estimates. Selecting the beds, path attachment, symmetry and
stationary tower on one retained history is likewise an incoming
construction requirement. The identity proves the connection between
those specified objects; it does not select new ones.
